\documentclass[10pt]{article}

\usepackage[margin=0.65in]{geometry}
\usepackage{amsmath,amssymb,mathtools}
\usepackage{booktabs,tabularx}
\usepackage{microtype}
\usepackage{multicol}
\usepackage[hidelinks]{hyperref}

\setlength{\parindent}{0pt}
\setlength{\parskip}{2pt}
\setlength{\columnsep}{0.22in}
\setlength{\abovedisplayskip}{3pt}
\setlength{\belowdisplayskip}{3pt}
\setlength{\abovedisplayshortskip}{2pt}
\setlength{\belowdisplayshortskip}{2pt}

\newenvironment{tightitemize}{%
  \begin{itemize}
  \setlength{\itemsep}{0pt}%
  \setlength{\parskip}{0pt}%
  \setlength{\parsep}{0pt}%
  \setlength{\topsep}{0pt}%
  \setlength{\partopsep}{0pt}%
}{\end{itemize}}

\newcommand{\code}[1]{\texttt{#1}}
\newcommand{\E}{\mathbb{E}}
\newcommand{\1}{\mathbf{1}}

\title{\vspace{-1.2em}Run v51 (DRL-Swing-Options): Actor, Schedules, and the \texorpdfstring{$\E[Q_T]$}{E[Q_T]} Proxy\vspace{-0.6em}}
\author{\vspace{-1.2em}}
\date{\vspace{-1.6em}}

\begin{document}
\small
\maketitle
\vspace{-0.8em}

\begin{multicols}{2}

\textbf{Scope.} This note documents (i) the actor network architecture, (ii) the learning-rate and exploration-noise schedules, and (iii) the HHK-based approximation used for $\E[Q_T]$. Assumes \code{runv51.sh} hyperparameters and the current implementations in \code{src/networks.py}, \code{src/agent.py}, \code{src/swing\_env.py}, \code{run.py}.

\vspace{0.2em}
\textbf{Run v51 essentials.}
\begin{tightitemize}
  \item Contract: $K{=}1$, $T{=}0.0833$, $n{=}22$, $(q_{\min},q_{\max}){=}(0,2)$, $(Q_{\min},Q_{\max}){=}(0,20)$, $r{=}0.05$, no refraction, $c{=}0$ (no convex cost).
  \item Actor/Critic: width $64$, depth $2$, activation SiLU, normalization RMSNorm, action head \texttt{tanh01}.
  \item LR: $(\eta_a,\eta_c)=(1.6{\times}10^{-4},\,9.0{\times}10^{-5})$, warmup $E_w{=}1024$ episodes, cosine horizon $E_h{=}40000$, final fraction $f{=}0.20$, $\eta_{\min}{=}10^{-6}$.
  \item Exploration: pre-squash Gaussian noise with $\sigma_0{=}1.30$, floor $\sigma_{\min}{=}0.26$, plateau $P{=}3200$ episodes.
\end{tightitemize}

\section*{1.\;Actor network (\code{src/networks.py:145})}
\textbf{State and action spaces.} The environment observation contains (at least) the intrinsic value proxy
\[
s_t[0] = S_t - K
\quad\text{and also }S_t\text{ explicitly (see }\code{src/swing\_env.py}\text{).}
\]
The actor outputs a normalized exercise intensity $a_t\in[0,1]$, mapped to a contract quantity
\[
q_t = q_{\min} + a_t\,(q_{\max}-q_{\min})
\quad\Rightarrow\quad
\text{(v51)}~q_t = 2a_t.
\]

\textbf{Architecture (v51).} Let $d_s$ be the state dimension and $d_a{=}1$ the action dimension. The actor is a 2-layer MLP with RMSNorm and SiLU:
\[
h_1 = \operatorname{SiLU}(\operatorname{RMSNorm}(W_1 s + b_1)),\qquad
h_2 = \operatorname{SiLU}(\operatorname{RMSNorm}(W_2 h_1 + b_2)),
\]
\[
z = W_3 h_2 + b_3,\qquad
a^{\text{raw}} = \tfrac12(\tanh(z)+1)\in[0,1].
\]
\textbf{Initialization.} Hidden linear layers use orthogonal initialization (gain matched to activation), and the final linear layer is initialized uniformly in $[-3\cdot 10^{-3},\,3\cdot 10^{-3}]$ (small initial actions).

\textbf{Profitability gate (hard constraint + STE).} The actor applies a hard gate which (forward-pass) forces $a_t=0$ when the immediate profit is non-positive, but (backward-pass) uses a straight-through estimator (STE) so gradients flow as if ungated. With v51 settings ($c{=}0$, $q_{\min}\ge 0$, \texttt{tanh01} head), the gate simplifies to:
\[
a_t = a^{\text{raw}}_t\,\1\{S_t-K>0\}
\quad\text{(forward, exactly),}
\]
implemented as
\[
a_t = a^{\text{raw}}_t + \big(a^{\text{raw}}_t\1\{S_t-K>0\}-a^{\text{raw}}_t\big)_{\text{stop-grad}}.
\]
The environment also masks actions to zero when the realized net payoff is $\le 0$ (\code{src/swing\_env.py}), consistent with the gate when $c{=}0$.

\section*{2.\;LR decay (warmup + cosine) (\code{src/agent.py})}
In v51, actor and critic optimizers are AdamW with decoupled weight decay; the LR is controlled by a per-episode multiplicative factor $g(e)$ (scheduler step uses 1-based episode indexing):
\[
\eta(e)=\eta_0\,g(e),\qquad e=1,2,\dots
\]
with
\[
g(e)=
\begin{cases}
\dfrac{e}{E_w}, & e\le E_w,\\[0.6em]
f+(1-f)\,\dfrac{1+\cos\!\big(\pi\,\frac{e-E_w}{E_h-E_w}\big)}{2}, & E_w< e < E_h,\\[0.6em]
f, & e\ge E_h,
\end{cases}
\qquad
g(e)\leftarrow \max\!\Big(g(e),\,\dfrac{\eta_{\min}}{\eta_0}\Big).
\]
Numerically (v51): $E_w{=}1024$, $E_h{=}40000$, $f{=}0.20$, $\eta_{\min}{=}10^{-6}$, $(\eta_{0,a},\eta_{0,c})=(1.6{\times}10^{-4},\,9.0{\times}10^{-5})$. The floor is inactive since $\eta_{\min}/\eta_0 \ll f$ for both actor and critic.

\section*{3.\;Exploration noise decay (\code{src/agent.py:\_pre\_noise\_sigma})}
Noise is applied \emph{to the actor pre-activation} $z$ (before \texttt{tanh01} squashing and before the profitability gate):
\[
z \leftarrow z + \sigma(e)\,\xi,\qquad \xi\sim\mathcal{N}(0,I).
\]
The pre-squash standard deviation $\sigma(e)$ is:
\[
\sigma(e)=
\begin{cases}
\sigma_0, & e < P,\\[0.4em]
\sigma_{\min} + \dfrac{\sigma_0-\sigma_{\min}}{1 + \frac{e-P}{P}}, & e\ge P,
\end{cases}
\]
so the post-plateau decay is hyperbolic in $e$ (half-life of order $P$). In v51: $(\sigma_0,\sigma_{\min},P)=(1.30,0.26,3200)$, hence $\sigma(6400)\approx 0.26+\frac{1.04}{2}=0.78$.

\section*{4.\;Approximating \texorpdfstring{$\E[Q_T]$}{E[Q\_T]} (\code{src/swing\_env.py:47})}
\textbf{Objective.} $Q_T:=\sum_{k=1}^n q_k$ is the cumulative exercised volume. Under v51 constraints ($Q_{\min}{=}0$, no refraction, linear payoff/cost), the policy is close to ``bang--bang'': exercise $q_{\max}$ on in-the-money (ITM) dates until the global cap $Q_{\max}$ is reached. The helper \code{approximate\_Q\_T} returns a fast proxy for $\E[Q_T]$ used to initialize the actor head mean in \code{run.py:warmup\_calibrate\_actor\_outputs}.

\textbf{HHK model input.} HHK is used in the form
\[
S_t=\exp(f(t)+X_t+Y_t),\qquad Z_t:=X_t+Y_t,
\]
where $X$ is Gaussian OU and $Y$ is a mean-reverting jump component (see parameters in \code{runv51.sh}: $\alpha{=}12$, $\sigma{=}1.2$, $\beta{=}150$, $\lambda{=}6$, $\mu_J{=}0.3$, $S_0{=}1$).

\textbf{Step A: marginal ITM probabilities.} For each exercise date $t_k$ (implemented on a grid $t_k=kT/n$), compute
\[
p_k := \mathbb{P}(S_{t_k}>K)=\mathbb{P}\!\big(Z_{t_k}>\log K - f(t_k)\big).
\]
If $\lambda\le 0$ (no jumps), $Z_{t_k}$ is Gaussian and $p_k$ is a normal tail probability. Otherwise, the code uses a Lugannani--Rice saddlepoint approximation: solve $K'( \hat\theta)=x$ for $x=\log K-f(t_k)$, where $K(\theta)=\log \E[e^{\theta Z_{t_k}}]$, then approximate $\mathbb{P}(Z_{t_k}\le x)$ by
\[
\Phi(w)+\phi(w)\Big(\frac1w-\frac1u\Big),\quad
w=\operatorname{sign}(\hat\theta)\sqrt{2(\hat\theta x-K(\hat\theta))},\quad
u=\hat\theta\sqrt{K''(\hat\theta)},
\]
and return $p_k\approx 1-\text{CDF}$.

\textbf{Step B: dependence proxy via exchangeable indicators.} Let $I_k=\1\{S_{t_k}>K\}$. The sum $N=\sum_{k=1}^n I_k$ is approximated by a Beta--Binomial distribution with:
\[
\bar p=\frac1n\sum_{k=1}^n p_k,
\qquad \rho_I \approx \text{avg}_{i<j}\operatorname{Corr}(I_i,I_j).
\]
The mapping $\operatorname{Corr}(Z_i,Z_j)\mapsto \operatorname{Corr}(I_i,I_j)$ is done via a Gaussian copula at thresholds $\tau_k=\Phi^{-1}(1-p_k)$ using a bivariate normal CDF (sampling up to \texttt{max\_pairs=4096} pairs).
Then Beta--Binomial parameters are moment-matched:
\[
s=\frac{1-\rho_I}{\rho_I},\qquad a=\bar p\,s,\qquad b=(1-\bar p)\,s,\qquad N\sim \text{BetaBin}(n,a,b).
\]

\textbf{Step C: truncated count-to-volume map.} Define $m=\lfloor Q_{\max}/q_{\max}\rfloor$ and $q_{\text{last}}=Q_{\max}-m q_{\max}$. Then
\[
\E[Q_T] \approx q_{\max}\,\E[\min(m,N)] + q_{\text{last}}\;\mathbb{P}(N>m),
\]
computed by summing the Beta--Binomial pmf over $k=0,\dots,n$ (and finally clipped to $[0,Q_{\max}]$). In v51, $q_{\max}{=}2$, $Q_{\max}{=}20$ so $m{=}10$ and $q_{\text{last}}{=}0$, i.e.
\[
\E[Q_T] \approx 2\,\E[\min(10,N)].
\]
\textbf{Usage in initialization.} \code{run.py} sets a target per-step volume $\E[Q_T]/n$ and calibrates the actor head so that the average normalized action is $\approx \text{normalize}(\E[Q_T]/n)$; for v51 this is $(\E[Q_T]/n)/2$.

\end{multicols}
\end{document}
